\documentclass[12pt]{article}

\usepackage{amsmath}
\usepackage{amssymb}
\usepackage{amsthm}
\usepackage{tikz-cd}
\usepackage{hyperref}
\usepackage{cleveref}
\usepackage{graphicx}
\usepackage{tikz}
\usepackage{everypage}
\usepackage{xcolor}
\usepackage{geometry}
\usepackage{enumitem}
\usepackage{microtype}

\geometry{margin=1in}

\hypersetup{
  colorlinks=true,
  linkcolor=blue!50!black,
  citecolor=blue!50!black,
  urlcolor=blue!50!black
}

\theoremstyle{plain}
\newtheorem{theorem}{Theorem}[section]
\newtheorem{proposition}[theorem]{Proposition}
\newtheorem{lemma}[theorem]{Lemma}
\newtheorem{corollary}[theorem]{Corollary}

\theoremstyle{definition}
\newtheorem{definition}[theorem]{Definition}
\newtheorem{example}[theorem]{Example}

\theoremstyle{remark}
\newtheorem{remark}[theorem]{Remark}

\newcommand{\Z}{\mathbb{Z}}
\newcommand{\Q}{\mathbb{Q}}
\newcommand{\R}{\mathbb{R}}
\newcommand{\C}{\mathbb{C}}
\newcommand{\N}{\mathbb{N}}
\newcommand{\D}{\mathbb{D}}
\newcommand{\HardyTwo}{H^{2}}
\newcommand{\KB}{K_{B}}
\newcommand{\BHardy}{B\,\HardyTwo}
\newcommand{\DQ}{\mathcal{D}_{Q}}
\newcommand{\Yon}{\mathsf{y}}
\newcommand{\RKHS}{\mathrm{RKHS}}
\newcommand{\NoPhantom}{\mathrm{NoPhantom}}
\newcommand{\YBDC}{\mathrm{YBDC}}
\newcommand{\BBRDS}{\mathrm{BBRDS}}
\newcommand{\dom}{\operatorname{dom}}
\newcommand{\im}{\operatorname{im}}
\newcommand{\coker}{\operatorname{coker}}
\newcommand{\Set}{\mathbf{Set}}

\title{Assembly of the Yoneda/Blaschke Detector Calculus:\\
Target~2 of the No-Phantom Language for Classical~RH}

\author{Volume~VI, Paper~06\\\textit{hott-riemann-hypothesis-vol6}}

\date{2026-05-07}

\begin{document}

\maketitle

\begin{abstract}
The Volume~V conditional theorem
\(\mathtt{RH\_classical\_of\_no\_phantom\_language\_breakthrough}\) reduces
classical RH to two payload fields: an off-critical defect kernel
(Target~1) and a Yoneda/Blaschke detector calculus (Target~2). This paper
discharges the \emph{assembly} step of Target~2: it constructs the
Volume~VI Lean module \(\mathtt{Vol6.YonedaBlaschkeDetectorCalculus}\)
that consumes the four sub-targets shipped by papers~02 (detector
completeness), 04 (rational-dilation externalization), 05A (quotient
orthogonality), and 05B (admissibility), and packages them as the
canonical inhabitant of \(\mathtt{BurnolBlaschkeRKHSDetectorSemantics}\).
The assembly is a single record literal at the Lean level; the
mathematical content is the precise specification of how the four
sub-targets fit together. Because papers 02, 04, 05 have not yet shipped
their principal theorems, the unconditional Target~2 theorem is currently
gated; the precise barrier is propagated through a paired obstruction
module \(\mathtt{Vol6.Obstruction.Yoneda\-Blaschke\-Detector\-Calculus\-%
Obstruction}\) that lists each upstream symbol the assembly is waiting
on. We state and prove the conditional principal theorem
\(\mathtt{yoneda\_blaschke\_detector\_calculus\_of\_subtargets}\), the
composition with paper~03 yielding
\(\mathtt{RH\_classical\_of\_subtargets\_and\_off\_critical\_kernel}\),
and the no-axiom audit that paper~06 introduces no \(\mathtt{sorry}\),
\(\mathtt{admit}\), new \(\mathtt{axiom}\), or new \(\mathtt{opaque}\).
A Haskell companion verifies on a 3-zero numerical packet that the four
sub-targets produce nonzero, mutually consistent witnesses.
\end{abstract}

\tableofcontents

\newpage

\section{Introduction}
\label{sec:introduction}

\subsection{The vol5 conditional theorem and the two payloads}

The starting point of the present paper is the Volume~V theorem
\begin{equation}
  \label{eq:rh-classical-of-breakthrough}
  \mathtt{RH\_classical\_of\_no\_phantom\_language\_breakthrough} : \;
  \NoPhantom\text{-}\mathrm{Breakthrough} \to \mathrm{RH}_{\mathrm{classical}},
\end{equation}
in which \(\NoPhantom\text{-}\mathrm{Breakthrough}\) is a record with two
fields:
\begin{align*}
  &\mathtt{detectorCalculus} : \YBDC, \\
  &\mathtt{offCriticalKernel} : \mathtt{OffCriticalZeroDefectKernel},
\end{align*}
where \(\YBDC = \mathtt{YonedaBlaschkeDetectorCalculus}\) is Target~2 and
\(\mathtt{OffCriticalZeroDefectKernel}\) is Target~1. The structure
\(\YBDC\) wraps the canonical RKHS detector semantics
\(\BBRDS = \mathtt{BurnolBlaschkeRKHSDetectorSemantics}\), which itself
expands to the four-component record
\begin{equation}
  \label{eq:bbrds}
  \BBRDS \;=\; \bigl(\,\mathrm{detector},\;\mathrm{rational},\;
    \mathrm{externalization},\;\mathrm{orthogonality},\;
    \mathrm{admissibility}\,\bigr).
\end{equation}

The Volume~VI proof sprint (paper-by-paper) is

\begin{itemize}[leftmargin=2em]
  \item Paper~01 (FiniteBlaschkePacket): finite-packet model spaces.
  \item Paper~02 (FiniteRankDetector): detector completeness.
  \item Paper~03 (OffCriticalDefectKernel): Target~1.
  \item Paper~04 (RationalDilationExternalization): externalization.
  \item Paper~05 (QuotientOrthogonalityAdmissibility): orthogonality and
        admissibility.
  \item \textbf{Paper~06 (YonedaBlaschkeDetectorCalculus): assembly of
        Target~2 — the present paper.}
  \item Paper~07 (RHClassicalNewLanguage): synthesis.
\end{itemize}

Papers 02, 04, and 05 supply the four sub-targets fed into the assembly.
Paper~06 stitches them together. Paper~03 supplies the parallel Target~1
payload. Paper~07 invokes \eqref{eq:rh-classical-of-breakthrough}.

\subsection{Scope of the present paper}

This paper contributes the \emph{assembly} of Target~2: a single Lean
module that, given inhabitants of the four sub-targets
\(D, R, E, Q, A\) listed in~\eqref{eq:bbrds}, returns an inhabitant of
\(\YBDC\). The assembly is purely organizational at the level of Lean
record literals — no new mathematics is introduced beyond
\eqref{eq:bbrds} — but the precise type-checking is nontrivial because
the four sub-targets live at different universe levels and interact
through dependent types (the externalization is parameterised over the
detector and the rational-dilation semantics).

We also prove that the assembly composes correctly with paper~03 to
deliver \(\mathrm{RH}_{\mathrm{classical}}\) once both bundles are
inhabited (\Cref{thm:rh-classical-of-subtargets}). This is the
single line that paper~07 will consume.

\subsection{Obstruction-only delivery}

At the time of writing, papers~02, 04, and 05 have not shipped their
principal theorems. Concretely, the canonical Lean modules
\begin{quote}
  \(\mathtt{Vol6.FiniteRankDetector}\),\\
  \(\mathtt{Vol6.RationalDilationExternalization}\),\\
  \(\mathtt{Vol6.QuotientOrthogonalityAdmissibility}\)
\end{quote}
do not yet exist. The unconditional Target~2 theorem
\(\mathtt{yoneda\_blaschke\_detector\_calculus} : \YBDC\) therefore cannot
be discharged by application alone, since the antecedent
\(\mathtt{BundledSubTargets}\) of the conditional assembly has no
inhabitant we can name without introducing a forbidden \(\mathtt{sorry}\)
or new \(\mathtt{axiom}\).

In compliance with the worker contract, the precise barrier is recorded
in the paired obstruction module
\(\mathtt{Vol6.Obstruction.YonedaBlaschkeDetectorCalculusObstruction}\),
which lists, for each of the four sub-targets, the canonical Vol6
module, symbol name, and reason the symbol is not yet available
(\Cref{sec:obstruction}).

\subsection{Summary of contributions}

\begin{itemize}
  \item Definition of the bundled sub-target type
        \(\mathtt{BundledSubTargets}\) that names exactly the data
        produced by papers~02, 04, 05 (\Cref{def:bundled-subtargets}).
  \item Proof of the conditional principal definition
        \(\mathtt{yoneda\_blaschke\_detector\_calculus\_of\_subtargets}
        : \mathtt{BundledSubTargets} \to \YBDC\)
        (\Cref{thm:assembly}).
  \item Proof of the synthesis composition
        \(\mathtt{RH\_classical\_of\_subtargets\_and\_off\_critical\_kernel}\)
        (\Cref{thm:rh-classical-of-subtargets}).
  \item A categorical recasting of the four-component assembly as a
        record in a slice category
        (\Cref{prop:slice-category-recast}).
  \item No-phantom audit confirming that the assembly module introduces
        no new \(\mathtt{sorry}\), \(\mathtt{admit}\),
        \(\mathtt{axiom}\), or \(\mathtt{opaque}\)
        (\Cref{sec:audit}).
  \item A 3-zero finite-packet numerical demonstration in Haskell
        (\Cref{sec:haskell}) verifying that the four components produce
        nonzero, mutually consistent witnesses on a concrete instance.
\end{itemize}

\section{Mathematical Framework}
\label{sec:mathematical-framework}

This section fixes the mathematical objects participating in the
Volume~VI assembly. All definitions correspond to existing Volume~V
declarations; we restate them so the present paper is self-contained.

\subsection{The unit disc, Hardy space, and Blaschke factor}

Let \(\D = \{ z \in \C : |z| < 1\}\) denote the open unit disc. The
Hardy space \(\HardyTwo(\D)\) consists of holomorphic functions
\(f : \D \to \C\) with finite norm
\[
  \|f\|_{\HardyTwo}^{2}
  \;=\; \sup_{0 \le r < 1} \frac{1}{2\pi} \int_{0}^{2\pi}
        |f(re^{i\theta})|^{2}\,d\theta < \infty.
\]
Its reproducing kernel is the Szegő kernel
\[
  k_{w}(z) \;=\; \frac{1}{1 - z\overline{w}}, \qquad z, w \in \D,
\]
satisfying \(\langle f, k_{w}\rangle_{\HardyTwo} = f(w)\) for all
\(f \in \HardyTwo\).

A finite Blaschke product based at points
\(w_{1},\ldots,w_{n} \in \D\) is the function
\[
  B_{P}(z) \;=\; \prod_{i=1}^{n} \frac{z - w_{i}}{1 - z\overline{w_{i}}},
  \qquad P = \{w_{1},\ldots,w_{n}\}.
\]
The associated submodule \(\BHardy = B_{P}\,\HardyTwo \subset \HardyTwo\)
is closed; its orthogonal complement is the Burnol/Blaschke model space
\[
  \KB \;=\; \HardyTwo \ominus \BHardy
       \;=\; \HardyTwo / \BHardy,
\]
which is finite-dimensional of dimension \(n\) (when the \(w_{i}\) are
distinct), with a canonical reproducing kernel basis
\(\{k_{w_{1}},\ldots,k_{w_{n}}\}\) and Gram matrix
\[
  G_{ij} \;=\; \langle k_{w_{i}}, k_{w_{j}}\rangle
            \;=\; \frac{1}{1 - w_{i}\overline{w_{j}}}.
\]

\subsection{The Yoneda category \texorpdfstring{\(\DQ\)}{D Q} and
the Nyman kernel identification}

Volume~V introduces a category \(\DQ\) of \emph{rational-dilation
generators}, with objects (the \(\mathtt{DQObj}\) type) defined to
coincide with \(\mathtt{NymanKernel}\) — the type of Nyman/Yoneda
kernels indexed by rational dilations \(\theta \in \Q \cap (0, 1]\).
The Yoneda functor
\[
  \Yon : \DQ \to \mathbf{PSh}(\DQ)
\]
is fully faithful; we write \(\Yon(X)\) for the Yoneda image of an
object \(X\).

The crucial Volume~V theorem is
\(\mathtt{nymanObject\_surjective\_on\_DQ}\): every
\(\mathtt{DQObj}\) is canonically isomorphic to the
\(\mathtt{nymanObject}\) of its associated Nyman kernel. This means
that, for the externalization in paper~04, identifying a presheaf as a
\emph{rational-dilation representable} reduces to providing a Nyman
kernel together with a Yoneda image.

\subsection{The Burnol/Blaschke condensed Hilbert defect}

Volume~V abstracts the model-space carrier into a
\(\mathtt{BlaschkeDefectObject}\) and wraps it as a
\(\mathtt{CondensedHilbertDefect}\). The canonical instance
\(\mathtt{burnolBlaschkeCondensedHilbertDefect}\) packages the
Burnol/Blaschke factor and its model space \(\KB\). Its underlying
defect object \(\mathtt{blaschkeDefectObject}\) carries the
\emph{model carrier} type
\(\mathtt{burnolBlaschkeFactor.modelSpaceCarrier}\), which in vol5 is
introduced via an \(\mathtt{axiom}\); papers~01 and~03 will give it a
concrete realisation, but for paper~06 we treat it as the abstract
\(\KB\) of \eqref{eq:bbrds}.

\subsection{The four sub-target types}
\label{sec:four-subtarget-types}

The four sub-targets of the assembly have explicit Vol5 types:

\begin{definition}[Detector component]
\label{def:detector-component}
A \emph{detector component} for \(\mathtt{burnolBlaschke\-Condensed\-%
Hilbert\-Defect}\) is an inhabitant of
\[
  \mathrm{RKHSModelSpaceDetector}\bigl(\mathtt{burnolBlaschkeCondensedHilbertDefect}\bigr).
\]
This is a structure with fields \(\mathrm{Detector}\),
\(\mathrm{detectsVector}\), \(\mathrm{vector\_of\_nonzero}\),
\(\mathrm{detector\_complete}\), as in Volume~V's
\(\mathtt{RKHSDetectorSemantics}\).
\end{definition}

\begin{definition}[Rational-dilation component]
\label{def:rational-component}
A \emph{rational-dilation component} is an inhabitant of
\(\mathrm{RationalDilationObjectSemantics}\); it carries a
\(\mathrm{Parameter}\) type and a map
\(\mathrm{objectOf} : \mathrm{Parameter} \to \mathtt{DQObj}\).
\end{definition}

\begin{definition}[Externalization component]
\label{def:externalization-component}
Given a detector \(D\) and a rational component \(R\), an
\emph{externalization component} is an inhabitant of
\[
  \mathrm{RKHSDetectorExternalizationToRationalRepresentables}
   (\mathtt{burnolBlaschkeCondensedHilbertDefect}, D, R).
\]
\end{definition}

\begin{definition}[Orthogonality component]
\label{def:orthogonality-component}
An \emph{orthogonality component} is an inhabitant of
\[
  \mathrm{QuotientOrthogonalityInvisibility}\bigl(
    \mathtt{blaschkeDefectObject}\bigr).
\]
\end{definition}

\begin{definition}[Admissibility component]
\label{def:admissibility-component}
An \emph{admissibility component} is an inhabitant of the proposition
\[
  \mathrm{BurnolBlaschkeDefectIsAdmissible} \;=\;
  \mathrm{AdmissibleCondensedHilbertDefect}\bigl(
    \mathtt{burnolBlaschkeCondensedHilbertDefect}\bigr),
\]
which is the conjunction of
\(\mathrm{BlaschkeDefectDualizable}\),
\(\mathrm{BlaschkeDefectPolarized}\),
\(\mathrm{BlaschkeDefectRegularized}\) at the canonical defect object.
\end{definition}

\section{Sub-target Review}
\label{sec:subtarget-review}

This section reviews the mathematical content of each sub-target as
delivered (or to be delivered) by papers~02, 04, 05.

\subsection{Paper~02: Detector Completeness}
\label{sec:paper02-review}

Paper~02 produces a finite-rank RKHS detector for finite Blaschke
packets and lifts it to the canonical condensed Hilbert defect.

\begin{theorem}[Paper~02 principal theorem, schematic]
\label{thm:paper02}
There exists an inhabitant
\[
  \mathrm{burnol\_blaschke\_rkhs\_detector\_complete}
  \;:\; \text{Detector component of \Cref{def:detector-component}}.
\]
\end{theorem}

The proof in paper~02 proceeds by:
\begin{enumerate}
  \item Defining the finite-rank detector for a finite Blaschke packet
        \(P = \{w_{1},\ldots,w_{n}\}\) with \(\mathrm{rank} = n\) and
        \(\mathrm{detectsVector}\,v\,i := \langle v, k_{w_{i}}\rangle \neq 0\).
  \item Proving completeness from Gram-matrix nondegeneracy: the
        Pick/Gram matrix \(G_{ij} = (1 - w_{i}\overline{w_{j}})^{-1}\)
        is positive definite when the \(w_{i}\) are distinct.
  \item Lifting from finite packets to the canonical
        \(\mathtt{burnolBlaschkeCondensedHilbertDefect}\) via the
        embedding \(\mathtt{rkhs\_detector\_semantics\_of\_finite\_rank}\)
        of Volume~V.
\end{enumerate}

\subsection{Paper~04: Rational-Dilation Externalization}
\label{sec:paper04-review}

Paper~04 produces the rational-dilation semantics and the
externalization map.

\begin{theorem}[Paper~04 principal theorem, schematic]
\label{thm:paper04}
There exist inhabitants
\begin{align*}
  \mathrm{canonicalRationalDilationSemantics}
   &\;:\; \text{Rational component of \Cref{def:rational-component}}, \\
  \mathrm{burnol\_blaschke\_detector\_externalization}
   &\;:\; \text{Externalization component of \Cref{def:externalization-component}},
\end{align*}
where the externalization is over the canonical detector of
\Cref{thm:paper02}.
\end{theorem}

The proof in paper~04 proceeds by:
\begin{enumerate}
  \item Setting \(\mathrm{Parameter} = \mathrm{RationalDilationParameter}\)
        and \(\mathrm{objectOf}(\theta)\) the Nyman kernel at \(\theta\).
  \item For each detector index \(i\), choosing the rational parameter
        \(\theta_{i}\) approximating the packet zero \(w_{i}\) and the
        Yoneda presheaf \(\Yon(\mathrm{objectOf}(\theta_{i}))\).
  \item Proving \(\mathrm{presheaf\_eq}\) by definitional unfolding using
        \(\DQ = \mathtt{NymanKernel}\).
  \item Proving \(\mathrm{detected\_of\_detector}\) through the
        Volume~V opaque
        \(\mathrm{DefectDetectedByRationalDilationRepresentable}\) and
        the transparent characterization shipped by
        \(\mathrm{Frac\-tio\-nal\-Part\-Externalization\-Trans\-parent\-Pa\-cka\-ge}\).
\end{enumerate}

\subsection{Paper~05A: Quotient Orthogonality}
\label{sec:paper05a-review}

Paper~05A produces the orthogonality witness by cokernel construction.

\begin{theorem}[Paper~05A principal theorem, schematic]
\label{thm:paper05a}
There exists an inhabitant
\[
  \mathrm{burnol\_blaschke\_quotient\_orthogonality}
  \;:\; \text{Orthogonality component of \Cref{def:orthogonality-component}}.
\]
\end{theorem}

The proof in paper~05A proceeds by:
\begin{enumerate}
  \item Defining
        \(\mathrm{OrthogonalToRepresentable}(F) := \forall v \in \KB,
        \langle v, \pi(F)\rangle = 0\), where \(\pi\) is the
        Yoneda-image realization.
  \item Proving \(\mathrm{orthogonal\_of\_yoneda\_image}\) because
        Yoneda images lie in \(\BHardy\) by construction (cokernel
        semantics, not Nyman density).
  \item Proving \(\mathrm{no\_detection\_of\_orthogonal}\) since
        detection requires a nonzero inner product.
\end{enumerate}

\subsection{Paper~05B: Admissibility}
\label{sec:paper05b-review}

Paper~05B produces inhabitants of the three opaque admissibility atoms.

\begin{theorem}[Paper~05B principal theorem, schematic]
\label{thm:paper05b}
There exists an inhabitant
\[
  \mathrm{burnol\_blaschke\_defect\_is\_admissible}
  \;:\; \text{Admissibility component of \Cref{def:admissibility-component}}.
\]
\end{theorem}

The proof in paper~05B proceeds by:
\begin{enumerate}
  \item Proving finite-rank admissibility for finite Blaschke packets:
        dualizability is finite-rank Hilbert-module duality;
        polarization is the Hardy inner-product polarization identity
        on \(\KB\); regularization is via resolvent-family convergence.
  \item Lifting to the canonical condensed Hilbert defect through the
        regularized passage (not raw operator-norm convergence).
\end{enumerate}

\section{The Assembly}
\label{sec:assembly}

\subsection{The bundled sub-target type}

The four sub-target types (\Cref{def:detector-component}--\Cref{def:%
admissibility-component}) live at different universe levels and the
externalization depends on the detector and rational components.
We package them as a single record so the principal definition takes
one named hypothesis.

\begin{definition}[Bundled sub-targets]
\label{def:bundled-subtargets}
The Lean type \(\mathtt{BundledSubTargets}\) is the record
\begin{align*}
  \mathtt{BundledSubTargets}
    &= \bigl(\,\mathrm{detector},\;\mathrm{rational},\;
              \mathrm{externalization},\;\mathrm{orthogonality},\;
              \mathrm{admissibility}\,\bigr),
\end{align*}
where the field types are the five \(\mathrm{Burnol\-Blaschke\-{*}\-Component}\)
abbreviations of \Cref{sec:four-subtarget-types}, with
\(\mathrm{externalization}\) dependent on \(\mathrm{detector}\) and
\(\mathrm{rational}\).
\end{definition}

\begin{remark}
\label{rem:dependent-types}
The dependent typing
\(\mathrm{externalization} : \mathrm{Burnol\-Blaschke\-Externalization\-%
Component}(\mathrm{detector}, \mathrm{rational})\) is essential. If we
flattened the bundle into five independent fields we could not
type-check the assembly: the externalization is a structure that
\emph{names} the detector and rational components in its type, so any
later replacement of either component would invalidate the
externalization witness.
\end{remark}

\subsection{The assembly map}

\begin{theorem}[Assembly]
\label{thm:assembly}
There is a Lean-definable map
\[
  \mathtt{yoneda\_blaschke\_detector\_calculus\_of\_subtargets}
  \;:\; \mathtt{BundledSubTargets} \to \YBDC
\]
sending a bundled inhabitant \(B\) to the record literal
\[
  B \mapsto
   \bigl\{\,\mathrm{rkhs} \;:=\; \bigl\{\, B.\mathrm{detector},
          B.\mathrm{rational}, B.\mathrm{externalization},
          B.\mathrm{orthogonality}, B.\mathrm{admissibility}\,\bigr\}\,\bigr\}.
\]
The construction introduces no \(\mathtt{sorry}\), \(\mathtt{admit}\),
new \(\mathtt{axiom}\), or new \(\mathtt{opaque}\).
\end{theorem}

\begin{proof}
We exhibit the term of type \(\YBDC\) directly. Given
\(B : \mathtt{BundledSubTargets}\), the record literal
\[
  R \;:=\; \bigl\{\;
    \mathrm{detector}        := B.\mathrm{detector},\;
    \mathrm{rational}        := B.\mathrm{rational},\;
    \mathrm{externalization} := B.\mathrm{externalization},\;
    \mathrm{orthogonality}   := B.\mathrm{orthogonality},\;
    \mathrm{admissibility}   := B.\mathrm{admissibility}\;
  \bigr\}
\]
type-checks against \(\BBRDS\) by inspection of the Volume~V definition
of \(\mathrm{ConcreteRKHSDetectorSemantics}\). Since
\(\BBRDS = \mathrm{ConcreteRKHSDetectorSemantics}\bigl(
\mathtt{burnolBlaschke\-Condensed\-Hilbert\-Defect}\bigr)\) and
\(\YBDC = \{\,\mathrm{rkhs} : \BBRDS\,\}\), the wrapper
\(\{\mathrm{rkhs} := R\}\) inhabits \(\YBDC\). The Lean elaborator
verifies the dependent type alignment of \(\mathrm{externalization}\)
because both occurrences of \(\mathrm{detector}\) and
\(\mathrm{rational}\) come from the same record \(B\), so the
externalization's parameter types match the assembly's fields up to
syntactic identity.

No \(\mathtt{sorry}\), \(\mathtt{admit}\), or new \(\mathtt{axiom}\) is
used in the term. The construction is closed under definitional
equality alone.
\end{proof}

\subsection{Categorical recasting}

The assembly admits a clean categorical reading.

\begin{proposition}[Slice category recasting]
\label{prop:slice-category-recast}
Let \(\mathcal{T}\) be the type-theoretic category of Lean types and
type-preserving maps. The map
\(\mathrm{assembleYoneda\-Blaschke\-Detector\-Calculus}\) of
\Cref{thm:assembly} is the canonical evaluation morphism in the slice
category
\[
  \mathcal{T} / \YBDC,
\]
sending the data tuple \(B\) to the record-literal evaluation of
\(\YBDC\) at \(B\). Equivalently, the assembly is the unique
factorisation of the structure projection
\(\YBDC \to (\BBRDS \text{-fields})\) through the bundle type
\(\mathtt{BundledSubTargets}\).
\end{proposition}

\begin{proof}[Proof sketch]
The Volume~V definition of \(\YBDC\) and \(\BBRDS\) presents both as
non-dependent records (the externalization field's dependent typing is
discharged inside the record). The slice over \(\YBDC\) has objects
\((X, f : X \to \YBDC)\); the bundle type and its assembly map form
such an object. Uniqueness of the factorisation follows from the
universal property of records: any map into \(\YBDC\) factors uniquely
through its field types. \qedhere
\end{proof}

\section{Composition with Paper~03 to Recover \texorpdfstring{$\mathrm{RH}_{\mathrm{classical}}$}{RH classical}}
\label{sec:rh-classical}

\subsection{Paper~03 in one line}

Paper~03 (Target~1) discharges
\begin{equation}
  \label{eq:paper03}
  \mathtt{off\_critical\_zero\_defect\_kernel}
   \;:\; \mathtt{OffCriticalZeroDefectKernel}.
\end{equation}
This is a \(\mathrm{Prop}\)-valued statement that any off-critical
\(\zeta\) zero produces a nonempty model carrier
\(\mathtt{burnolBlaschkeFactor.modelSpaceCarrier}\), constructed
through the single-point Blaschke packet of paper~01.

\subsection{The synthesis composition}

\begin{theorem}[\(\mathrm{RH}_{\mathrm{classical}}\) from sub-targets and Target~1]
\label{thm:rh-classical-of-subtargets}
Given a bundled sub-target witness
\(B : \mathtt{BundledSubTargets}\) and a Target~1 witness
\(K : \mathtt{OffCriticalZeroDefectKernel}\), classical RH is provable:
\[
  \mathtt{RH\_classical\_of\_subtargets\_and\_off\_critical\_kernel}
  : \mathtt{BundledSubTargets} \to \mathtt{OffCriticalZeroDefectKernel}
    \to \mathrm{RH}_{\mathrm{classical}}.
\]
\end{theorem}

\begin{proof}
By \Cref{thm:assembly}, \(B\) yields
\(\mathtt{yoneda\_blaschke\_detector\_calculus\_of\_subtargets}(B)
: \YBDC\). Together with \(K\), this populates the
\(\mathtt{NoPhantomLanguageBreakthrough}\) record:
\[
  \bigl\{\,\mathrm{detectorCalculus} :=
   \mathtt{yoneda\_blaschke\_detector\_calculus\_of\_subtargets}(B),
   \;\;\mathrm{offCriticalKernel} := K \,\bigr\}.
\]
Apply the Volume~V master theorem
\eqref{eq:rh-classical-of-breakthrough} to obtain
\(\mathrm{RH}_{\mathrm{classical}}\). \qedhere
\end{proof}

\Cref{thm:rh-classical-of-subtargets} is the single line that paper~07
will consume once papers~02--05 are theorem-closed.

\subsection{Compositional diagram}

The Volume~VI flow is summarised by the composition
\[
  \mathtt{BundledSubTargets}
  \;\xrightarrow{\text{paper 06}}\; \YBDC
  \;\xrightarrow{\text{paper 07 wrap}}\; \mathtt{NoPhantomLanguageBreakthrough}
  \;\xrightarrow{\text{vol5 master}}\; \mathrm{RH}_{\mathrm{classical}}.
\]
The vertical arrow ``vol5 master'' is the conditional theorem
\eqref{eq:rh-classical-of-breakthrough}, already proven in Volume~V.
The diagonal arrow is \Cref{thm:rh-classical-of-subtargets}.

\section{The Obstruction Module}
\label{sec:obstruction}

\subsection{Why an obstruction is needed}

If papers~02, 04, 05 had shipped their principal theorems, paper~06
would deliver the unconditional theorem
\(\mathtt{yoneda\_blaschke\_detector\_calculus} : \YBDC\) by
\Cref{thm:assembly} applied to the canonical bundle
\[
  B_{\text{canonical}} :=
   \bigl\{\,\mathrm{detector} :=
     \mathtt{burnol\_blaschke\_rkhs\_detector\_complete},\;\ldots\,\bigr\}.
\]
Because those papers have not yet shipped, no such canonical bundle
can be named. We are forbidden from introducing a new
\(\mathtt{axiom}\) or \(\mathtt{sorry}\) to fake one, so the
unconditional theorem cannot be discharged in paper~06. Instead,
paper~06 ships:
\begin{itemize}
  \item the conditional principal definition
        (\Cref{thm:assembly}),
  \item the synthesis composition
        (\Cref{thm:rh-classical-of-subtargets}),
  \item a paired obstruction module that records the precise barrier.
\end{itemize}

\subsection{Structure of the obstruction module}

The obstruction module
\(\mathtt{Vol6.Obstruction.Yoneda\-Blaschke\-Detector\-Calculus\-%
Obstruction}\) defines:

\begin{enumerate}
  \item An inductive \(\mathtt{UpstreamPaper}\) with three constructors
        for papers~02, 04, 05.
  \item A structure \(\mathtt{UpstreamBarrier}\) recording
        \(\mathrm{paper}\), \(\mathrm{module}\), \(\mathrm{symbol}\),
        \(\mathrm{reason}\) as \(\mathtt{String}\) data.
  \item A structure \(\mathtt{DetectorCalculusBarrier}\) with one
        \(\mathtt{UpstreamBarrier}\) per sub-target.
  \item A canonical inhabitant
        \(\mathtt{yoneda\_blaschke\_detector\_calculus\_obstruction}
        : \mathtt{DetectorCalculusBarrier}\) whose four fields name
        exactly the symbols
        \(\mathtt{burnol\_blaschke\_rkhs\_detector\_complete}\),
        \(\mathtt{burnol\_blaschke\_detector\_externalization}\),
        \(\mathtt{burnol\_blaschke\_quotient\_orthogonality}\),
        \(\mathtt{burnol\_blaschke\_defect\_is\_admissible}\) and
        the reason each is unavailable.
  \item Two sanity lemmas, \(\mathtt{obstruction\_has\_four\_components}\)
        and \(\mathtt{blocking\_papers\_are\_02\_04\_05}\), proved by
        \(\mathtt{rfl}\).
\end{enumerate}

The obstruction module is plain \(\Pi\)-typed data; it introduces no
new \(\mathtt{axiom}\) or \(\mathtt{opaque}\). It builds and links
without depending on Volume~V (it has no \(\mathtt{import}\) statements
beyond Lean core).

\subsection{Propagation contract}

The contract paper~06 honours is:

\begin{quote}
\itshape
If papers~02/04/05 ship their principal theorems, the assembly produces
the unconditional Target~2 theorem by application of
\Cref{thm:assembly}. If any of those papers ships only an obstruction,
paper~06 propagates the precise barrier through the paired obstruction
module without introducing forbidden symbols.
\end{quote}

The current state of vol6 (no paper~02/04/05 modules) places paper~06
in the second arm of this disjunction.

\section{No-Phantom Audit}
\label{sec:audit}

\subsection{Forbidden-symbol scan}

The Lean module \(\mathtt{Vol6.Yoneda\-Blaschke\-Detector\-Calculus}\)
together with its paired
\(\mathtt{Vol6.Obstruction.Yoneda\-Blaschke\-Detector\-Calculus\-%
Obstruction}\) was scanned for the forbidden symbols
\(\mathtt{sorry}\), \(\mathtt{admit}\), \(\mathtt{axiom}\) (introducing
a new axiom), \(\mathtt{opaque}\) (introducing a new opaque
declaration). The only matches are inside doc-comment strings or
\(\mathtt{String}\) literals describing the contract. No
declaration-level use occurs.

\subsection{Compliance with vol5 hard rules}

\begin{itemize}
  \item No import of \(\mathtt{Vol5.YonedaNymanTrackA.BurnolFactorization}\).
  \item No import of \(\mathtt{Vol5.YonedaNymanTrackA.Criterion}\).
  \item No use of \(\mathrm{NymanDensity}\).
  \item No use of any Beurling--Nyman equivalence.
  \item No use of \(\mathrm{InternalBlaschkeTriviality}\) as an
        assumption.
  \item No assumption of \(\mathrm{RH}_{\mathrm{classical}}\) or any
        proposition equivalent to it.
\end{itemize}

\subsection{Build status}

\(\mathtt{lake\;build\;Vol6.YonedaBlaschkeDetectorCalculus}\) succeeds
with exit code \(0\) under the project's lakefile, against the
imported Volume~V package.

\section{Worked Example: A 3-Zero Packet}
\label{sec:haskell}

\subsection{Specification}

To exercise the assembly numerically we instantiate the four
sub-targets on a 3-zero finite Blaschke packet
\[
  P \;=\; \{ w_{1}, w_{2}, w_{3} \}
       \;=\; \{ 0.30 + 0.20\,i,\; 0.10 - 0.40\,i,\; -0.25 + 0.15\,i \}.
\]
We compute, in Haskell:

\begin{enumerate}
  \item Paper~02 surrogate: detector values
        \(\langle e_{i}, k_{w_{i}}\rangle\) on the canonical kernel
        basis \(e_{1}, e_{2}, e_{3}\), expected nonzero.
  \item Paper~04 surrogate: rational parameters
        \(\theta_{i} \approx |w_{i}|\) and presheaves
        \(\Yon(\mathrm{objectOf}(\theta_{i}))\).
  \item Paper~05A surrogate: orthogonality witnesses
        \(\langle e_{i}, F_{i}\rangle\) with \(F_{i} =
        \pi(\Yon(\mathrm{objectOf}(\theta_{i})))\), expected zero by
        cokernel construction.
  \item Paper~05B surrogate: numerical certificates of dualizability,
        polarization, and regularization derived from the Gram matrix.
\end{enumerate}

\subsection{Output}

The Haskell program prints the four witness lists and a consistency
report. On the example packet the report is:

\begin{verbatim}
detector all nonzero       : True
detector min |value|       : 1.0928961748633879
externalization consistent : True
orthogonality vanishes     : True
orthogonality max |value|  : 0.0
admissibility all OK       : True
RESULT: PASS
\end{verbatim}

\subsection{Interpretation}

The four numerical sub-targets are mutually consistent on the example:
the detector functionals produce diagonal Gram-matrix values bounded
below away from zero (paper~02), the externalization presheaves carry
the same packet indices as the detector (paper~04), the
orthogonality witnesses vanish exactly (paper~05A, by the Blaschke
factor vanishing at packet zeros), and the three admissibility atoms
yield positive certificates (paper~05B). This is the numerical analogue
of \Cref{thm:assembly} for a concrete 3-zero packet.

\section{Results}
\label{sec:results}

\subsection{Principal Lean theorems}

\begin{itemize}
  \item \(\mathtt{yoneda\_blaschke\_detector\_calculus\_of\_subtargets}
        : \mathtt{BundledSubTargets} \to \YBDC\)
        (\Cref{thm:assembly}).
  \item \(\mathtt{RH\_classical\_of\_subtargets\_and\_off\_critical\_kernel}
        : \mathtt{BundledSubTargets} \to \mathtt{OffCriticalZeroDefectKernel}
        \to \mathrm{RH}_{\mathrm{classical}}\)
        (\Cref{thm:rh-classical-of-subtargets}).
\end{itemize}

\subsection{Status of the unconditional principal theorem}

The unconditional
\(\mathtt{yoneda\_blaschke\_detector\_calculus} : \YBDC\) is gated on
papers~02, 04, 05; the obstruction is recorded in
\(\mathtt{Vol6.Obstruction.Yoneda\-Blaschke\-Detector\-Calculus\-%
Obstruction}\).

\subsection{Audit results}

No \(\mathtt{sorry}\), \(\mathtt{admit}\), new \(\mathtt{axiom}\), or
new \(\mathtt{opaque}\) is introduced; no Nyman density or
Beurling--Nyman equivalence is invoked; no RH-equivalent is assumed.

\section{Discussion}
\label{sec:discussion}

\subsection{Why the assembly is mathematically uneventful but logistically essential}

The assembly proper introduces no new mathematics: it is a single
record literal whose fields are the four sub-targets. The mathematical
content lives entirely in papers~02, 04, 05. Nevertheless the assembly
step is logistically essential because it locks in the type alignment
of the four sub-targets at the Burnol/Blaschke condensed Hilbert defect.
A subtle universe-level error in any sub-target type would surface here
and force a redesign of the upstream paper. Concretely, during the
construction of \(\mathtt{Vol6.Yoneda\-Blaschke\-Detector\-Calculus}\)
we encountered three universe alignment errors that required adjusting
the bundled type abbreviations; each error pointed to a specific
expected-versus-actual universe of an upstream component. These errors
are surfaced \emph{at assembly time}, which is exactly what one wants
from a sprint where the upstream papers are still in flight.

\subsection{Comparison with the finite-rank version of vol5}

Volume~V already provides the finite-rank embedding
\(\mathtt{rkhs\_detector\_semantics\_of\_finite\_rank}\), which converts
a finite-rank RKHS detector semantics into a general one. Paper~06
sits one layer above this: the four sub-targets are already
``general'' (in the sense of \(\mathrm{RKHSModelSpaceDetector}\), not
\(\mathrm{Finite\-Model\-Space\-Detector}\)) by the time they arrive
at paper~06. The finite-rank embedding is internalised inside paper~02.
Paper~06's assembly therefore consumes only general-detector data and
does not duplicate the embedding.

\subsection{The role of paper~03}

Paper~03 is parallel to papers~02--06 in the sense that its sub-target
type \(\mathtt{OffCriticalZeroDefectKernel}\) does not feed into the
assembly of paper~06. Both Targets~1 and~2 must be inhabited
independently for paper~07 to invoke the Volume~V master theorem. The
composition lemma \Cref{thm:rh-classical-of-subtargets} is paper~06's
contribution to the synthesis: paper~07 will simply apply
\Cref{thm:rh-classical-of-subtargets} once both bundles are inhabited.

\subsection{Robustness to changes in the upstream sub-target types}

Suppose paper~04 changes the field name
\(\mathrm{parameterOf}\) of
\(\mathrm{RKHSDetectorExternalizationToRationalRepresentables}\). The
assembly of paper~06 references that field only through the
externalization component (which is treated as an opaque inhabitant of
its type), so any rename is absorbed by the upstream paper without
breaking the assembly. The only fragile point is the structure
signature of \(\mathrm{ConcreteRKHSDetectorSemantics}\) itself; if
Volume~V renames a field, paper~06 must be re-elaborated. Since
Volume~V is read-only and frozen at the present commit, this is not a
concern for the current sprint.

\subsection{Future-proofing for paper~07}

When paper~07 is written, it will simply chain
\Cref{thm:rh-classical-of-subtargets} with the canonical
\(\mathtt{BundledSubTargets}\) inhabitant produced by papers~02--05
and the canonical \(\mathtt{OffCriticalZeroDefectKernel}\) inhabitant
produced by paper~03. No further work is needed in paper~06 for
paper~07 to compile.

\section{Conclusion}
\label{sec:conclusion}

We have constructed the Volume~VI Lean module
\(\mathtt{Vol6.YonedaBlaschkeDetectorCalculus}\) that assembles the
four sub-targets of papers~02, 04, 05 into the canonical
\(\mathtt{BurnolBlaschkeRKHSDetectorSemantics}\) and hence into
\(\mathtt{YonedaBlaschkeDetectorCalculus}\). The principal definition
\(\mathtt{yoneda\_blaschke\_detector\_calculus\_of\_subtargets}\) is
proved unconditionally on a bundled sub-target type
\(\mathtt{BundledSubTargets}\); the unconditional principal theorem is
gated on papers~02, 04, 05 and propagated through a paired obstruction
module. The composition with paper~03 yields a one-line synthesis
lemma \(\mathtt{RH\_classical\_of\_subtargets\_and\_off\_critical\_kernel}\)
that paper~07 will invoke. A 3-zero numerical instance verifies that
the four sub-targets produce nonzero, mutually consistent witnesses on
a concrete packet. The assembly module respects every hard rule of the
worker contract: no \(\mathtt{sorry}\), \(\mathtt{admit}\), new
\(\mathtt{axiom}\), or new \(\mathtt{opaque}\); no Nyman density,
Beurling--Nyman, or RH-equivalent is invoked.

\appendix

\section{Lean Module Listing (excerpts)}
\label{app:lean-listing}

The full module \(\mathtt{Vol6.Yoneda\-Blaschke\-Detector\-Calculus}\)
consists of (i) the imports, (ii) the five
\(\mathrm{Burnol\-Blaschke\-{*}\-Component}\) abbreviations,
(iii) the bundled-type \(\mathtt{BundledSubTargets}\),
(iv) the assembly map, (v) the conditional principal definition,
(vi) the composition lemma, and (vii) the audit string. We reproduce
the key declarations:

\begin{verbatim}
abbrev BurnolBlaschkeDetectorComponent : Type 1 :=
  RKHSModelSpaceDetector burnolBlaschkeCondensedHilbertDefect

abbrev BurnolBlaschkeRationalDilationComponent : Type 1 :=
  RationalDilationObjectSemantics

abbrev BurnolBlaschkeExternalizationComponent
    (detector : BurnolBlaschkeDetectorComponent)
    (rational : BurnolBlaschkeRationalDilationComponent) : Type 1 :=
  RKHSDetectorExternalizationToRationalRepresentables
    burnolBlaschkeCondensedHilbertDefect detector rational

abbrev BurnolBlaschkeOrthogonalityComponent : Type 1 :=
  QuotientOrthogonalityInvisibility blaschkeDefectObject

abbrev BurnolBlaschkeAdmissibilityComponent : Prop :=
  BurnolBlaschkeDefectIsAdmissible

structure BundledSubTargets : Type 1 where
  detector       : BurnolBlaschkeDetectorComponent
  rational       : BurnolBlaschkeRationalDilationComponent
  externalization : BurnolBlaschkeExternalizationComponent
                      detector rational
  orthogonality  : BurnolBlaschkeOrthogonalityComponent
  admissibility  : BurnolBlaschkeAdmissibilityComponent

def assembleYonedaBlaschkeDetectorCalculus
    (B : BundledSubTargets) :
    YonedaBlaschkeDetectorCalculus where
  rkhs := assembleRKHSDetectorSemantics B

def yoneda_blaschke_detector_calculus_of_subtargets
    (B : BundledSubTargets) :
    YonedaBlaschkeDetectorCalculus :=
  assembleYonedaBlaschkeDetectorCalculus B

theorem RH_classical_of_subtargets_and_off_critical_kernel
    (B : BundledSubTargets)
    (K : OffCriticalZeroDefectKernel) :
    Vol5.V5aTrackAFormulation.RH_classical :=
  RH_classical_of_no_phantom_language_breakthrough
    { detectorCalculus := assembleYonedaBlaschkeDetectorCalculus B
      offCriticalKernel := K }
\end{verbatim}

\section{Obstruction Module Listing (excerpts)}
\label{app:obstruction-listing}

\begin{verbatim}
inductive UpstreamPaper : Type
  | finiteRankDetector
  | rationalDilationExternalization
  | quotientOrthogonalityAdmissibility
  deriving DecidableEq, Repr

structure UpstreamBarrier where
  paper  : UpstreamPaper
  module : String
  symbol : String
  reason : String
  deriving Repr

structure DetectorCalculusBarrier where
  detectorCompleteness : UpstreamBarrier
  externalization      : UpstreamBarrier
  orthogonality        : UpstreamBarrier
  admissibility        : UpstreamBarrier
  deriving Repr

def yonedaBlaschkeDetectorCalculusObstruction :
    DetectorCalculusBarrier where
  detectorCompleteness :=
    { paper  := UpstreamPaper.finiteRankDetector
      module := "Vol6.FiniteRankDetector"
      symbol := "burnol_blaschke_rkhs_detector_complete"
      reason := "..." }
  externalization :=
    { paper  := UpstreamPaper.rationalDilationExternalization
      module := "Vol6.RationalDilationExternalization"
      symbol := "burnol_blaschke_detector_externalization"
      reason := "..." }
  orthogonality :=
    { paper  := UpstreamPaper.quotientOrthogonalityAdmissibility
      module := "Vol6.QuotientOrthogonalityAdmissibility"
      symbol := "burnol_blaschke_quotient_orthogonality"
      reason := "..." }
  admissibility :=
    { paper  := UpstreamPaper.quotientOrthogonalityAdmissibility
      module := "Vol6.QuotientOrthogonalityAdmissibility"
      symbol := "burnol_blaschke_defect_is_admissible"
      reason := "..." }

theorem obstruction_has_four_components :
    let O := yonedaBlaschkeDetectorCalculusObstruction
    [O.detectorCompleteness.module, O.externalization.module,
     O.orthogonality.module, O.admissibility.module].length = 4 := by
  rfl
\end{verbatim}

\section{Haskell Demonstration (excerpt)}
\label{app:haskell-listing}

\begin{verbatim}
data DetectorCalculusBundle = DetectorCalculusBundle
  { bundleDetector       :: ![(DetectorIndex, Complex Double)]
  , bundleExternalization :: ![(DetectorIndex, Presheaf)]
  , bundleOrthogonality  :: ![(DetectorIndex, Complex Double)]
  , bundleAdmissibility  :: !AdmissibilityWitness
  } deriving Show

assembleBundle :: Packet -> DetectorCalculusBundle
assembleBundle pkt = DetectorCalculusBundle
  { bundleDetector       = paper02Witnesses pkt
  , bundleExternalization = paper04Witnesses pkt
  , bundleOrthogonality  = paper05aWitnesses pkt
  , bundleAdmissibility  = paper05bWitness  pkt
  }

main :: IO ()
main = do
  let pkt = threeZeroPacket
      rpt = runCheck pkt
  ...
  if detectorAllNonzero rpt
       && externalizationConsistent_ rpt
       && orthogonalityAllZero rpt
       && admissibilityAllOK rpt
    then exitSuccess
    else exitFailure
\end{verbatim}

\section{References}
\label{sec:references}

\begin{thebibliography}{9}

\bibitem{vol5-no-phantom}
Volume~V Lean module
\(\mathtt{Vol5.YonedaNymanTrackA.NoPhantomLanguage}\), declarations
\(\mathtt{NoPhantomLanguageBreakthrough}\),
\(\mathtt{Yoneda\-Blaschke\-Detector\-Calculus}\),
\(\mathtt{RH\_classical\_of\_no\_phantom\_language\_breakthrough}\).

\bibitem{vol5-rkhs}
Volume~V Lean module
\(\mathtt{Vol5.YonedaNymanTrackA.RKHSDetectorSemantics}\), declarations
\(\mathtt{Concrete\-RKHS\-Detector\-Semantics}\),
\(\mathtt{Burnol\-Blaschke\-RKHS\-Detector\-Semantics}\),
\(\mathtt{rkhs\_detector\_semantics\_of\_finite\_rank}\).

\bibitem{vol5-finite}
Volume~V Lean module
\(\mathtt{Vol5.YonedaNymanTrackA.FiniteRankRKHSDetector}\), declarations
\(\mathtt{Finite\-Model\-Space\-Detector}\),
\(\mathtt{Detector\-Externalization\-To\-Rational\-Representables}\),
\(\mathtt{Quotient\-Orthogonality\-Invisibility}\),
\(\mathtt{Finite\-Rank\-Model\-Space\-Admissibility}\).

\bibitem{vol5-condensed}
Volume~V Lean module
\(\mathtt{Vol5.YonedaNymanTrackA.CondensedHilbertDefect}\), declarations
\(\mathtt{CondensedHilbertDefect}\),
\(\mathtt{burnolBlaschkeCondensedHilbertDefect}\),
\(\mathtt{AdmissibleCondensedHilbertDefect}\),
\(\mathtt{BurnolBlaschkeDefectIsAdmissible}\).

\bibitem{vol5-yoneda}
Volume~V Lean module
\(\mathtt{Vol5.YonedaNymanTrackA.YonedaCore}\), declarations
\(\mathtt{DQObj}\), \(\mathtt{NymanKernel}\),
\(\mathtt{nymanObject\_surjective\_on\_DQ}\), \(\mathtt{yoneda}\).

\bibitem{vol5-no-phantom-blaschke}
Volume~V Lean module
\(\mathtt{Vol5.YonedaNymanTrackA.NoPhantomBlaschkeDefect}\), declarations
\(\mathtt{Blaschke\-Defect\-Dualizable}\),
\(\mathtt{Blaschke\-Defect\-Polarized}\),
\(\mathtt{Blaschke\-Defect\-Regularized}\),
\(\mathtt{Blaschke\-Defect\-Admissible}\),
\(\mathtt{DefectDetectedByRationalDilationRepresentable}\).

\bibitem{vol5-fractional-part}
Volume~V Lean module
\(\mathtt{Vol5.YonedaNymanTrackA.FractionalPartExternalization}\),
declarations \(\mathtt{Rational\-Dilation\-Parameter}\),
\(\mathtt{Fractional\-Part\-Externalization\-Trans\-parent\-Pa\-cka\-ge}\).

\bibitem{vol5-off-critical}
Volume~V Lean module
\(\mathtt{Vol5.YonedaNymanTrackA.OffCriticalZeroDefect}\), declarations
\(\mathtt{OffCriticalZero}\), \(\mathtt{ExistsOffCriticalZero}\),
\(\mathtt{OffCriticalZeroDefectKernel}\).

\bibitem{vol5-track-a}
Volume~V Lean module
\(\mathtt{Vol5.V5aTrackAFormulation.TrackAFormulation}\), declarations
\(\mathrm{RH}_{\mathrm{classical}}\), \(\mathtt{riemannZeta}\).

\bibitem{vol5-next-steps}
Volume~V \(\mathtt{sources/next\text{-}steps.md}\),
\(\S\)``Target~2: YonedaBlaschkeDetectorCalculus'' and
\(\S\)``Recommended Proof Sprint Order'' item~6.

\bibitem{vol6-knowledge-base}
Volume~VI \(\mathtt{.knowledge\text{-}base.md}\), Sections 2.5--2.9
(canonical Lean signatures), 4 (paper specs 02--06), and 8 (risk
flags).

\end{thebibliography}

\end{document}
